\section{Numerical implementation and accuracy checks}
\label{app:rewrite-numerical-details}

This section supplements the numerical method in
Appendix~\ref{app:rewrite-algorithm} of the main paper. I describe the
computational coordinates, the construction of numerical initial guesses,
and the checks used to accept the simulated trajectories.

\subsection{Computational coordinates and initial guesses}

For each $\sigma$, I first solve the existing-AI economy with efficiency
fixed at $B_0$ and research unavailable. A scalar root search in $\log K_0$
solves Equation~\eqref{eq:rewrite-rsi-initial-ratio} using the production
and monopoly equations. The resource constraint~\eqref{eq:rewrite-resource},
with $M=0$ and $\dot K=(n+\gamma)K$, gives pre-event consumption.
The post-activation problem inherits the stocks but selects consumption,
research, and the shadow price anew.

To represent positive variables and keep AI efficiency below its upper
bound, I use the computational coordinates
\[
 \left(\log K,\ \log\frac{B}{\overline B-B},\ \log C,\ \log q\right).
\]
The second coordinate expresses AI efficiency relative to the remaining
gap below its upper bound; it can grow without bound while AI efficiency
remains bounded. I subtract the long-run growth rate of aggregate capital
times $t$ from each coordinate: $n+\gamma$ in the labor-bottleneck regime
and $n+\mathcal R(\overline B)-\rho$ in the AI-dominated regime.
The efficiency coordinate has the same asymptotic slope because the
remaining gap $\overline B-B$ shrinks at that proportional rate.
Removing these trends makes the numerical coordinates converge to finite
limits without changing the economic variables or calendar time.

For each of the eight parameter configurations, I start near its
analytical limiting point and move the initial stocks to their prescribed
values through a sequence of auxiliary boundary-value problems, using
each solution as the initial guess for the next. This procedure changes
the initial conditions of the numerical problem, not the dates on the
economic trajectory.

For the competitive-to-monopoly comparison in
Section~\ref{subsec:rewrite-competitive-transition}, I instead obtain
$K_0$ from $r=\rho+\gamma$ with marginal-cost pricing $p_X=1/B_0$;
at $\sigma=1.5$, $K_0=6.535$. Pre-event consumption follows from the same
resource constraint. I then solve the monopoly BVP at these initial
stocks, using the RSI-activation trajectories only as numerical initial
guesses, not as restrictions on consumption or research. This procedure
also applies to the four-elasticity comparisons in the
\href{https://oliverpardo1979.github.io/ai-growth-and-labor/\#app:rewrite-competitive-transition}{digital edition}.

\subsection{Residuals and horizon checks}

To check whether the computed transition depends on the terminal date or
solver accuracy, I extend each computational horizon twice by 500 years and tighten the
collocation tolerance from $2\times10^{-6}$ to $10^{-8}$ and $10^{-9}$.
The final boundary tolerance is $10^{-11}$. All plotted values lie within
the solved interval; none is extrapolated beyond it.

I check the four dynamic equations independently of the solver's residual
report. Five-point finite differences approximate the time derivatives of
the four logarithmic coordinates above; I compare these derivatives with
the equations' right-hand sides expressed in the same coordinates.
The dynamic residual is the largest absolute discrepancy across
coordinates and checked dates, measured per year, with denser checks over
the first ten years. For the monopoly and research first-order conditions, I report the
absolute logarithmic deviations of $Bp_X(1-e_X)$ and
$q\chi\eta B^\eta M^{\eta-1}\psi(B)$ from one.
Acceptance requires dynamic residuals below $10^{-6}$ and both
first-order residuals below $10^{-9}$.

I also require the last horizon extension to change the four detrended
coordinates by less than $2\times10^{-5}$ over years 0--500, with the same
tolerance for changes in $\log C_0$ and $\log q_0$.
At the terminal date, the largest coordinate deviation from the
analytical limiting point must be below $10^{-4}$. This last check uses
the regime-specific normalized coordinates of
Appendix~\ref{app:rewrite-finite-existence}, with logarithms for coordinates
having positive limits and with $\tau=(AN)^{-1}$ scaled by
$\overline d/\overline B$. Here $\overline d$ is the limiting value of the
normalized efficiency gap $d$ in the corresponding regime. Unlike the
horizon comparison, this check measures proximity to the limiting
equilibrium, not agreement between two numerical solutions.
I also check the pre-event equilibrium conditions and the continuity of
production quantities and prices at RSI activation.

Small equation residuals do not establish that the developer prefers the
computed research path to every feasible alternative. I therefore check
feasibility and the sufficient conditions for global optimality in
Lemma~\ref{lem:rewrite-developer-verification}. The parameter condition
\eqref{eq:rewrite-research-optimality-condition} is verified in
Appendix~\ref{app:rewrite-numerical-support-checks} of the main paper for
all eight configurations.

To assess the behavior of the paths beyond the computational horizon, I check proximity
to the analytical limiting equilibrium at the terminal date and use the
asymptotic decay of both transversality
expressions at rate $n-\rho<0$. These checks support the numerical
approximation over an infinite horizon; they do not constitute a proof
of equilibrium existence from arbitrary initial stocks.
The competitive-to-monopoly paths satisfy the same residual,
horizon-refinement, transversality, and developer-optimality criteria.

\begin{table}[H]
\centering
\caption{Numerical accuracy of the illustrative RSI comparisons}
\label{tab:rewrite-rsi-illustrative-accuracy}
\small
\begin{tabular}{rrrrr}
\toprule
$\chi$ & $\sigma$ & Final horizon & Dynamic residual & Horizon change\\
\midrule
7.5 & 0.90 & 5,428.0 & $1.18\times10^{-9}$ & $3.06\times10^{-8}$\\
7.5 & 1.00 & 4,944.5 & $1.38\times10^{-9}$ & $3.46\times10^{-8}$\\
7.5 & 1.10 & 4,828.8 & $1.66\times10^{-9}$ & $3.73\times10^{-8}$\\
7.5 & 1.50 & 3,013.0 & $1.69\times10^{-9}$ & $1.56\times10^{-8}$\\
1.5 & 0.90 & 6,047.0 & $1.02\times10^{-9}$ & $4.30\times10^{-8}$\\
1.5 & 1.00 & 5,563.5 & $1.11\times10^{-9}$ & $3.80\times10^{-8}$\\
1.5 & 1.10 & 5,447.8 & $1.19\times10^{-9}$ & $4.08\times10^{-8}$\\
1.5 & 1.50 & 3,632.0 & $1.23\times10^{-9}$ & $1.80\times10^{-8}$\\
\bottomrule
\end{tabular}
\par\smallskip
\begin{minipage}{0.98\textwidth}\footnotesize
Horizons are measured in years. Dynamic residuals are the largest absolute
errors in the derivatives of the four logarithmic coordinates, across
the full-horizon and dense first-decade checks. Horizon change compares the
last two solutions in detrended logarithmic coordinates over years 0--500.
Reported error maxima are rounded upward. All eight paths pass the equation,
event-continuity, developer-optimality and long-run-continuation checks.
The $\sigma=1.50$ cases use the global Hamiltonian-support test and its
analytical long-run bound; the other cases pass the concavity test.
\end{minipage}
\end{table}

